\begin{claim}{F.27}{Proved}{All finite affine marginal images on an exact two-parameter character locus}
Use the actual negative free basis $X_1,\ldots,X_7$ of
F.24 and the based comparison $\mathcal B_+:V_-\to V_+$
of F.23.13. Put $Y_i^-=X_i$ and
$Y_i^+=\mathcal B_+(X_i)$. For every finite field $k$
and continuous $\rho:V_\pm\to\operatorname{Aff}_1(k)$,
let $\chi$ be its scalar character. The character locus
\[
 \chi(Y_1^\pm)=\chi(Y_4^\pm)=\chi(Y_5^\pm)
       =\chi(Y_6^\pm)=1,\qquad
 \chi(Y_2^\pm)\chi(Y_3^\pm)\chi(Y_7^\pm)=1
                                                               \tag{F.27.1}
\]
has the exact image equality
\[
                         \rho(\mathcal K_\pm)=\rho(R_\pm).
                                                               \tag{F.27.2}
\]
There is no restriction on the characteristic, extension
degree, or size of the original coefficient field.

More precisely, write a negative character in this locus as
$(1,s,t,1,1,1,c^{-1})$, where $c=st$, and put
$k_0=\mathbb F_p(s,t)$, $k_1=\mathbb F_p(c)$.
For $c\ne1$ and $k_1\subsetneq k_0$, let
$q_1=|k_1|$, $m=[k_0:k_1]$ and
$\chi_j=\chi^{q_1^j}$, $0\leq j<m$.
Writing $K=\mathcal K_-$ and $V=V_-$, the actual map
\[
 \bigoplus_{j=0}^{m-1}H^1_{\rm cont}(V,k_{0,\chi_j})
   \longrightarrow H^1_{\rm cont}(K,k_{0,\chi|K}),\qquad
                 ([d_j])_j\longmapsto\sum_j[d_j|K]
                                                               \tag{F.27.3}
\]
is injective. At the cocycle level, the kernel of
the summed restriction is exactly
\[
 \left\{\bigl(b_j(\chi_j-1)\bigr)_j:
                  b_j\in k_0,\ \sum_jb_j=0\right\}.
                                                               \tag{F.27.4}
\]
Thus its rank is $6m+1$ on $7m$ cocycle coordinates.
The proof uses the literal eleven F.23 source
automorphisms and specified source conjugates of
one of its eight negative seeds.
\end{claim}

\begin{proof}
\emph{The actual scalar fields and all exceptional cases.}
Work first with the negative groups and write
$K\leq R\leq V$. The four deleted generators
$X_4,\ldots,X_7$ normally generate $R$.
The normal generation and actual marginal maps
of F.21--F.22, as used in F.25.5 and F.26.5, give
\[
 \begin{gathered}
 C=\chi(K)=\chi(R)=\langle c^{-1}\rangle,\qquad
 D=\chi(V)=\langle s,t\rangle,\\
 k_1=\mathbb F_p(C)=\mathbb F_p(c),\qquad
 k_0=\mathbb F_p(D)=k_1(s)=k_1(t),\qquad c=st.
 \end{gathered}                                                \tag{F.27.5}
\]
Indeed the last equalities follow from $t=c/s$ and
$s=c/t$. When $c=1$, the $R$ image consists of
translations and is abelian; the actual normal
generation therefore gives (F.27.2). For $c\ne1$
and $k_0=k_1$, F.25 gives the same equality.
It remains to treat $c\ne1$ and $k_1\subsetneq k_0$.
Here $s,t\notin k_1$, by (F.27.5), and consequently
\[
             s\,t\,(s-1)(t-1)(t+1)(c-1)\ne0.
                                                               \tag{F.27.6}
\]
The factors $s,t$ are nonzero character values.
The elements $1,-1$ belong to $k_1$, so neither
can equal $s$ or $t$. This includes characteristic
two, where $t+1=t-1$. The case $c=1$ has already
been settled, so every factor in (F.27.6) is justified.

\emph{The exact coefficient map and principal quotient.}
All universal word calculations in this proof are
performed on the discrete free group on the seven
specified generators, with its canonical dense map
to $V=\widehat F_7$. The coefficient homomorphism
from the F.24 Laurent ring is
\[
 \begin{aligned}
 \mathbb Z[t_1^{\pm1},\ldots,t_7^{\pm1}]&\longrightarrow
                  A=\mathbb Z[s^{\pm1},t^{\pm1}],\\
 (t_1,t_2,t_3,t_4,t_5,t_6,t_7)
           &\longmapsto(1,s,t,1,1,1,(st)^{-1}).
 \end{aligned}                                                \tag{F.27.7}
\]
Thus a monomial with exponent vector $(a_i)$ maps
to $s^{a_2-a_7}t^{a_3-a_7}$. The universal
derivation on free words has $D(X_i)=e_i$ and
$D(uv)=D(u)+\chi(u)D(v)$. Evaluation of its
coefficients in a finite field gives the actual
cocycle value on that word. Every seven-tuple
of generator values defines a homomorphism from
the discrete free group to the finite affine
group and extends uniquely to its full profinite
completion. Consequently the continuous cocycle
space is exactly $k_0^7$. No continuous map to
an infinite Laurent ring is being used.

Put $\delta=(1,s,t,1,1,1,c^{-1})^{\mathsf T}
-\mathbf1$. Its last entry is
$\delta_7=(1-c)/c\ne0$.
The quotient by principal cocycles has the explicit
projection and section
\[
 \begin{aligned}
 P(v)&=\bigl(v_r-\delta_r\delta_7^{-1}v_7\bigr)_{r=1}^6,
       &S(w)&=(w_1,\ldots,w_6,0)^{\mathsf T},\\
 PS&=I_6,
       &SP(v)&=v-\delta\delta_7^{-1}v_7.
 \end{aligned}                                                \tag{F.27.8}
\]
These formulas identify
$Q_\chi=H^1_{\rm cont}(V,k_{0,\chi})$ with $k_0^6$.
Universally they are defined over $A[(1-st)^{-1}]$;
all finite-field evaluations here respect that
localization by (F.27.6).

Let $B_l(v)=u_l\alpha_l(v)u_l^{-1}$, $0\leq l\leq10$,
be the actual old-source automorphisms of F.23.
They preserve $K$, including under their inverses,
and fix the character because their action on
the abelianization of $V$ is the identity.
Let $A_l$ be their F.24 precomposition matrices
after (F.27.7), so the $i$th row is $D(B_l(X_i))$.
Then $A_l\delta=\delta$ and their actual cohomology
actions are
\[
                           \overline A_l=PA_lS.
                                                               \tag{F.27.9}
\]
The source indices are exactly
$0=\texttt{old0}$,
$1,\ldots,5=\texttt{old\_1\_0},\ldots,\texttt{old\_1\_4}$
and $6,\ldots,10=\texttt{old\_-1\_0},\ldots,
\texttt{old\_-1\_4}$.

\emph{A basis with explicit inverse and its actual source action.}
For $P_0=X_1\cdots X_7$, the character is one,
and its derivative row is $(1,1,s,c,c,c,c)$.
Every pure $\alpha_l$ fixes this word. Hence
evaluation on $P_0$ induces the cohomology
covector $b=(1,1,s,c,c,c)$, with
$b\overline A_l=\chi(u_l)b$.
For each $r\in\{1,4,5,6\}$, the pure Artin
action writes $\alpha_l(X_r)=a_{lr}X_ra_{lr}^{-1}$.
Since $\chi(X_r)=1$, its actual conjugated source
action gives
$d(B_l(X_r))=\chi(u_la_{lr})d(X_r)$.
Thus $e_1^*,e_4^*,e_5^*,e_6^*$ are four more
actual quotient covectors. Their common kernel
together with $b$ is the line generated by
\[
                      v_0=(0,1,-s^{-1},0,0,0)^{\mathsf T}.
                                                               \tag{F.27.10}
\]
Use the basis matrix $J$ with columns
\[
       v_0,\quad e_2,\quad e_5-ce_2,\quad e_4-ce_2,
                       \quad e_1-e_2,\quad e_6-ce_2.
                                                               \tag{F.27.11}
\]
Its inverse is the explicit map
\[
 J^{-1}(w)=(-sw_3,\,bw,\,w_5,\,w_4,\,w_1,\,w_6)^{\mathsf T}.
                                                               \tag{F.27.12}
\]
For completeness, $J(z_0,\ldots,z_5)^{\mathsf T}$ has
coordinates
\[
 \bigl(z_4,\ z_0+z_1-cz_2-cz_3-z_4-cz_5,
          \ -s^{-1}z_0,\ z_3,\ z_2,\ z_5\bigr)^{\mathsf T}.
                                                               \tag{F.27.13}
\]
Substitution in either direction gives the identity,
so these are actual mutually inverse coefficient maps.

The five displayed quotient covectors give the
following form for the source actions in this basis:
\[
 T_l=J^{-1}\overline A_lJ
       =\begin{pmatrix}\nu_l&r_l\\0&\Lambda_l\end{pmatrix},
 \qquad \Lambda_l\text{ diagonal},\qquad
 (\nu_0,\ldots,\nu_{10})=(1,s,t,1,1,1,s,t,1,1,1).
                                                               \tag{F.27.14}
\]
The last character is computed from the literal
source words. More explicitly, for the seven-coordinate
lift $\widetilde v_0=S(v_0)$ the verified identities are
\[
 \begin{gathered}
 A_l\widetilde v_0=\nu_l\widetilde v_0+\pi_l\delta,\\
 \pi_1=-c,\qquad \pi_2=\pi_7=ct,\qquad
 \pi_6=s(c-t-1),\qquad
 \pi_l=0\quad(l\notin\{1,2,6,7\}).
 \end{gathered}                                                \tag{F.27.15}
\]
The certificate records every entry of all eleven
$A_l$, $\overline A_l$ and $T_l$, with the coefficient
maps (F.27.7)--(F.27.13).

\emph{An integral determinant forcing the essential line.}
Take the additive matrix commutators
$E_j=T_1T_j-T_jT_1$ for $j=2,4,5,6,7$.
They kill the first coordinate line and have image
in that line, since the lower diagonal blocks
in (F.27.14) commute. Their first-row entries on
the other five coordinates, in that order, form
the exact matrix
\[
 \begin{gathered}
 gH,\qquad g=(s-1)(t-1),\\[-2pt]
 H=\begin{pmatrix}
 0&0&0&-(t+1)/t&st\\
 0&s&-1/t&0&0\\
 0&-1/t&0&st+1&0\\
 (s-1)/(st^2)&s(s-1)&-(s-1)/t&s-1&0\\
 (t+1)/(st^2)&0&0&0&st
 \end{pmatrix}.
 \end{gathered}                                                \tag{F.27.16}
\]
Direct determinant expansion gives the integral
Laurent identity
\[
 \det H=\frac{(s-1)(t-1)(t+1)}{t^4},\qquad
 \det(gH)=\frac{(s-1)^6(t-1)^6(t+1)}{t^4}.
                                                               \tag{F.27.17}
\]
All its factors are nonzero after every finite-field
evaluation under consideration, by (F.27.6).
The identities have integer coefficients and only
signed Laurent monomial denominators; no prime
is removed from the calculation.

Let $L_\chi=k_0v_0\leq Q_\chi$. It is invariant
by (F.27.14). Every nonzero subspace $W\leq Q_\chi$
invariant under the eleven actual source operators
contains $L_\chi$. Indeed, in the $J$ coordinates
the common kernel of the five $E_j$ is precisely
the first coordinate line, by (F.27.17).
If $W$ is contained in this line, it equals it.
Otherwise choose a vector of $W$ with nonzero
last five coordinates. Some $E_j$ sends it to
a nonzero vector of the first coordinate line;
this vector remains in $W$. Since the line has
dimension one, it is contained in $W$.
This proves the required essential-socle statement
by explicit source endomorphisms.

\emph{Actual seed evaluations and all Frobenius factors.}
Let $d_0$ be the continuous cocycle with generator
values $\widetilde v_0$. Take the literal negative
seed $a=u_{\texttt{cross\_4\_-1\_1}}$ of index $25$.
The actual seed calculation gives
\[
                   \chi(a)=c^{-1},\qquad d_0(a)=c-1.
                                                               \tag{F.27.18}
\]
Both are computed by the same decorated braid
words as in F.23; all eight seed rows are audited.
Since $a\in K$ and $B_1$ preserves $K$, put
\[
          \kappa=B_1(a)a^{-1}\in K,
          \qquad \kappa_r=B_1^r(\kappa)\in K\quad(r\geq0).
                                                               \tag{F.27.19}
\]
Their characters are one. Using $\pi_1=-c$ in
(F.27.15) gives
$d_0(B_1(a))=(s+1)(c-1)$, hence
$d_0(\kappa)=s(c-1)$.
For every character-one element $h$, (F.27.15)
also gives $d_0(B_1(h))=s\,d_0(h)$: the principal
correction evaluates to zero there. Induction yields
\[
                         d_0(\kappa_r)=s^{r+1}(c-1).
                                                               \tag{F.27.20}
\]
In row coordinates the first two evaluations
are $D(a)(A_1-I_7)$ and $D(a)(A_1-I_7)A_1$;
the independent audit verifies their character-one
principal vanishing and the stated values.

Put $s_j=s^{q_1^j}$ and $t_j=t^{q_1^j}$.
Since $c\in k_1$, $s_jt_j=c$ for every $j$.
The field equality $k_0=k_1(s)$ proves that
$s_0,\ldots,s_{m-1}$ are distinct. Indeed an
equality would make a nonidentity power of the
$k_1$-Frobenius fix $s$, and hence its generated
field $k_0$; the powers before $m$ are distinct
automorphisms of the field of $q_1^m$ elements.

The coefficient map on cocycles and on cohomology is
\[
                  d\longmapsto d^{q_1^j},
 \qquad k_{0,\chi}\longrightarrow k_{0,\chi_j}.
                                                               \tag{F.27.21}
\]
It is additive and semilinear, preserves the
cocycle rule, carries principals to principals,
and preserves continuity because the field is finite.
Its inverse is the power $q_1^{m-j}$, with the
identity used at $j=0$; the two composites are
the identity on $k_0$. It commutes with the actual
source precomposition operators. Thus the explicit
bases, inverse maps, commutator identities and
essential lines transport to every $\chi_j$.
Write $Q_j=H^1_{\rm cont}(V,k_{0,\chi_j})$ and
$L_j=k_0[d_j]\leq Q_j$, where $d_j=d_0^{q_1^j}$.
Equation (F.27.20) becomes
\[
              d_j(\kappa_r)=a_js_j^r,
                     \qquad a_j=s_j(c-1)\ne0.
                                                               \tag{F.27.22}
\]

The characters $\chi_j$ all agree on $K$, since
$\chi(K)=C\subset k_1^\times$. Evaluation at
$\kappa_0,\ldots,\kappa_{m-1}$ therefore defines
a map on the common target cohomology space:
principals vanish on each character-one element.
On the summed restrictions of the $m$ essential
lines its matrix is
\[
                  \bigl(a_js_j^r\bigr)_{
                        0\leq r,j<m},\qquad
 \det=\left(\prod_j a_j\right)
                   \prod_{j<l}(s_l-s_j)\ne0.
                                                               \tag{F.27.23}
\]
The coefficient inverse is explicit. Write
\[
 \prod_{l\ne j}\frac{X-s_l}{s_j-s_l}
       =\sum_{r=0}^{m-1}h_{jr}X^r.
 \quad\text{Then}\quad
       z_j=a_j^{-1}\sum_{r=0}^{m-1}h_{jr}y_r
                                                               \tag{F.27.24}
\]
is the inverse to
$y_r=\sum_j a_js_j^r z_j$.
Evaluation of the displayed interpolation polynomial
at every $s_l$ proves both inverse identities.
Every denominator in (F.27.24) is nonzero by
the proved distinctness and (F.27.22).

\emph{Summed restriction, the trace kernel and the affine image.}
The map (F.27.3) commutes with all eleven source
operators: on both source and target they are
the actual maps $[d]\mapsto[d\circ B_l]$.
Consequently its kernel is an invariant subspace
of the finite-dimensional direct sum $\bigoplus_jQ_j$.
Suppose that kernel were nonzero and choose a
nonzero invariant subspace $T$ inside it of
smallest dimension. Each coordinate projection
$T\to Q_j$ has invariant kernel and image.
Minimality makes every nonzero projection
injective, with an image having no proper
nonzero invariant subspace. The essential-line
result puts $L_j$ inside that image; therefore
the image equals $L_j$. All projections of $T$
thus lie in their respective $L_j$, and
$T\leq\bigoplus_jL_j$. But (F.27.23)--(F.27.24)
make evaluation after summed restriction injective
on this entire direct sum of lines. This
contradicts $T$ lying in the kernel and proves
(F.27.3).

Zero summed restriction of cocycles now implies
that each $d_j$ is principal. Their restrictions
have the common character $\chi|K$, so their
sum is $(\chi|K-1)\sum_jb_j$. Since $C\ne1$,
this vanishes exactly when $\sum_jb_j=0$.
This proves (F.27.4), including both directions.
The seven-coordinate description of each
cocycle space gives rank $7m-(m-1)=6m+1$.

For the trace map of F.26.2, a cocycle $d$ in
its kernel yields the tuple $(d^{q_1^j})_j$,
whose summed restriction is exactly
$\operatorname{Tr}_{k_0/k_1}(d|K)=0$.
Thus $d=b(\chi-1)$ by the result just proved.
Trace vanishing on $K$, together with
$\chi(K)=C\ne1$, gives $\operatorname{Tr}(b)=0$.
Conversely these principal cocycles have zero
trace restriction. Hence the precise trace
kernel is $\mathscr P_\chi$ of F.26.3.

Every proper affine marginal image in the current
locus would, by the explicit conjugation and
coefficient homomorphism of F.26.1--12, give a
nonprincipal element of that trace kernel over
its scalar field $k_0$. Those maps retain the
original scalar character, so all the displayed
locus identities and fields remain the same.
The proved trace kernel excludes such an element.
This establishes (F.27.2) in the negative case,
for every original finite coefficient field $k$.

\emph{Certificate and transport.}
The symbolic producer applies exactly (F.27.7)
to the fixed, previously audited F.24 universal
rows. The separate verifier imports neither
that producer nor a computer algebra package.
It evaluates all eleven actual old-source words
and eight actual negative seed words by generic
affine multiplication over the integral Laurent
ring $A$. It checks the inverse basis maps,
the principal-quotient action after clearing
$\delta_7$, all source and socle identities,
the actual seed evaluations, and (F.27.17)
by its $120$ signed permutation terms.
The verifier performs $699$ recorded checks.
Thus the determinant certificate is an exact
integer identity and applies in every characteristic;
its nonzero finite-field factors are separately
proved in (F.27.6). The evidence manifest binds
both programs, their outputs and stdout, the
original braid input and all accepted dependencies.

Finally, a positive representation in (F.27.1)
precomposes with the actual continuous isomorphism
$\mathcal B_+:V_-\to V_+$ to give exactly the
negative character just treated. By F.23.13,
this isomorphism and its inverse carry both
$\mathcal K_-$ and $R_-$ onto their positive
counterparts. The equality of the two actual
images therefore transports through these
specified maps, proving (F.27.2) for both signs.
\end{proof}
